% ═══ Results: Reconstruction Assessment and Negative Results ═══
\section{Reconstruction Assessment and Negative Results}
\label{sec:res-reconstruction-negative}

\subsection{Reconstruction is not achievable at current signal strength}
We assessed whether the circuits supply sufficient correct spatial constraints
for template-free structure determination by measuring recall (fraction of
native contacts recovered) and long-range precision.  Circuit recall@L = 0.028
versus published DCA's 0.218; median long-range precision@L/5 is 0.000 for
circuits versus 0.683 for PSICOV predictions.  Against the threshold of
median LR prec@L/5 $\geq 0.30$ and recall@L $\geq 0.10$ derived from
Jones et al.\ and the CONFOLD literature, the circuits do not meet the
criterion while published DCA does.

\subsection{Tension mechanism does not improve prediction}
An attention-style tension score, normalising each position's mutual
information by its strongest partner, did not improve circuit precision
(tension-augmented 0.154 vs.\ non-tension 0.165).  High-tension positions were
anti-correlated with structural hubs (median $\rho = -0.06$), indicating that
entropic attention concentrates on variable surface clusters rather than
buried cores.

\subsection{Iterative DI is unreliable on conserved-rich families}
Our mean-field implementation of iterative direct information produced
anti-predictive rankings (AUC $\approx 0.41$) on families with many fully
conserved columns, due to pseudo-inverse conditioning when variance approaches
zero for locked positions.  APC-corrected Frobenius norm was adopted as the
primary DCA score instead; this choice was validated against native contacts
and matches published mean-field benchmarks.

\subsection{Cross-dataset pilot}
A pilot evaluation on three EVmutation proteins produced inconclusive results
because correct parsing of the A2M alignment format requires handling
insert-state columns that do not map linearly to PDB residue numbers.
Structures were successfully fetched from RCSB and MI AUCs were positive for
two of three targets, but precision could not be meaningfully evaluated
without proper reference-frame mapping.  Full evaluation is deferred to
future work.
